% =============================================================================
% Valorisation des swaps OIS EUR (ESTR) : theorie et implementation
% Compiler avec :  pdflatex pricing_ois.tex  (deux passes)
% Les chiffres, tableaux et donnees des figures sont generes par
% generate_data.py a partir du code du depot (voir README en fin de document).
% =============================================================================
\documentclass[11pt,a4paper]{article}

\usepackage[utf8]{inputenc}
\usepackage[T1]{fontenc}
\usepackage{lmodern}
\usepackage[french]{babel}
\usepackage{textcomp}
\usepackage{amsmath,amssymb,amsthm,mathtools}
\usepackage{microtype}
\usepackage[a4paper,margin=2.5cm]{geometry}
\usepackage{booktabs,array,tabularx,longtable}
\usepackage{graphicx}
\usepackage{float}
\usepackage[dvipsnames]{xcolor}
\usepackage{enumitem}
\usepackage{listings}
\usepackage{tikz}
\usetikzlibrary{arrows.meta,positioning,shapes.geometric,fit,calc,babel}
\usepackage{pgfplots}
\pgfplotsset{compat=newest}
\pgfkeys{/pgf/number format/.cd, use comma}
\usepackage{fancyhdr}
\usepackage[font=small,labelfont=bf]{caption}
\usepackage[colorlinks=true,linkcolor=NavyBlue,urlcolor=NavyBlue,citecolor=NavyBlue]{hyperref}

% --- Mise en page ------------------------------------------------------------
\setlength{\parindent}{0pt}
\setlength{\parskip}{0.55em}
\setlength{\emergencystretch}{3em}
\setlist{itemsep=0.15em,topsep=0.3em}
\pagestyle{fancy}
\fancyhf{}
\fancyhead[L]{\small Valorisation des swaps OIS \texteuro STR}
\fancyhead[R]{\small\leftmark}
\fancyfoot[C]{\small\thepage}
\renewcommand{\headrulewidth}{0.3pt}

% --- Macros ------------------------------------------------------------------
\newcommand{\estr}{\texorpdfstring{\texteuro STR}{EUR STR}}
\newcommand{\code}[1]{\texttt{\detokenize{#1}}}
\newcommand{\pct}[1]{#1\,\%}
\newcommand{\eur}{\ifmmode\,\text{\texteuro}\else\,\texteuro\fi}
\newcommand{\DF}[1]{P(#1)}
\DeclareMathOperator{\adj}{adj}
\DeclareMathOperator{\adv}{adv}

\theoremstyle{plain}
\newtheorem{proposition}{Proposition}
\theoremstyle{definition}
\newtheorem{definition}{Définition}
\theoremstyle{remark}
\newtheorem*{remarque}{Remarque}

% --- Listings ----------------------------------------------------------------
\definecolor{codebg}{gray}{0.96}
\definecolor{codekw}{rgb}{0.10,0.25,0.55}
\definecolor{codecm}{rgb}{0.30,0.45,0.30}
\lstset{
  language=Python,
  basicstyle=\fontsize{7.6}{9.2}\selectfont\ttfamily,
  keywordstyle=\color{codekw}\bfseries,
  commentstyle=\color{codecm}\itshape,
  stringstyle=\color{BrickRed},
  backgroundcolor=\color{codebg},
  frame=single,
  rulecolor=\color{black!25},
  breaklines=true,
  columns=fullflexible,
  keepspaces=true,
  showstringspaces=false,
  aboveskip=0.6em,
  belowskip=0.4em,
  captionpos=b,
}
\renewcommand{\lstlistingname}{Listing}
\lstdefinelanguage{json}{morestring=[b]", sensitive=false}

% --- Figures -----------------------------------------------------------------
\pgfplotsset{
  every axis/.append style={
    grid=major, grid style={black!12}, tick label style={font=\footnotesize},
    label style={font=\footnotesize}, legend style={font=\footnotesize, draw=black!30},
  },
}
\tikzset{
  layer/.style={draw, rounded corners=3pt, minimum width=5.4cm, minimum height=0.9cm,
    align=center, fill=blue!6},
  proc/.style={draw, rounded corners=3pt, minimum width=2.9cm, minimum height=0.95cm,
    align=center, fill=blue!6, font=\small},
  data/.style={draw, rounded corners=3pt, minimum height=0.8cm, align=center,
    fill=orange!12, font=\small},
  flow/.style={-{Stealth[length=2.4mm]}, thick},
}

% --- Chiffres generes par generate_data.py -----------------------------------
% Generated by docs/latex/generate_data.py: do not edit by hand.
\newcommand{\MaxGapLogLinear}{$2{,}4\times10^{-13}$}
\newcommand{\MaxGapLinearZero}{$2{,}4\times10^{-12}$}
\newcommand{\PillarCount}{24}
\newcommand{\MaxIterations}{9}
\newcommand{\MaxResidual}{$6{,}0\times10^{-15}$}
\newcommand{\FiveYPV}{1\,902{,}54}
\newcommand{\FiveYFixedLeg}{$-$1\,070\,177{,}48}
\newcommand{\FiveYOvernightLeg}{1\,072\,080{,}01}
\newcommand{\FiveYParRate}{2{,}2540}
\newcommand{\FiveYAnnuity}{47\,563\,443{,}37}
\newcommand{\FiveYPVBP}{4\,756{,}34}
\newcommand{\FiveYStart}{2026-10-21}
\newcommand{\FiveYEnd}{2031-10-21}
\newcommand{\FiveYCashflowCount}{10}
\newcommand{\FiveYPVCheck}{1\,902{,}54}
\newcommand{\TwoYPV}{$-$1\,581{,}22}
\newcommand{\TwoYParRate}{1{,}9920}
\newcommand{\TwoYAnnuity}{19\,765\,217{,}95}
\newcommand{\TelStart}{2026-10-21}
\newcommand{\TelEnd}{2027-10-21}
\newcommand{\TelDays}{257}
\newcommand{\TelProduct}{1{,}01911180555556}
\newcommand{\TelRatio}{1{,}01911180555556}
\newcommand{\TelDiff}{$1{,}3\times10^{-15}$}
\newcommand{\ValuationDate}{2026-10-19}
\newcommand{\FixingCount}{203}
\newcommand{\QuoteCount}{24}


\title{\textbf{Valorisation des swaps OIS \estr{}}\\[0.4em]
  \large Théorie du pricing et implémentation dans le pricer de dérivés de taux}
\author{Projet \texttt{rate-derivatives-pricer} -- phase 1}
\date{\today}

\begin{document}
\maketitle
\thispagestyle{empty}

\begin{abstract}
\noindent
Ce document présente, d'une part, la théorie de la valorisation d'un swap de taux indexé sur un taux
overnight (OIS) en monnaie euro : capitalisation du taux \estr{}, valorisation par actualisation sur
une courbe unique, taux de par, annuité, construction de la courbe par \emph{bootstrap}. Il expose,
d'autre part, les choix de conception et l'implémentation de la phase~1 du projet : bibliothèque
quantitative en Python et NumPy, service FastAPI, client Python typé et tableau de bord Streamlit.
Chaque résultat chiffré (courbe, exemples, écarts) est produit par le code du dépôt, et
l'implémentation est validée contre QuantLib avec des écarts de l'ordre de \MaxGapLogLinear{} sur les
facteurs d'actualisation.

\medskip\noindent
\textbf{Avertissement.} Les données de marché utilisées (snapshot \code{EUR_2026-10-19}) sont
\emph{synthétiques} : aucun niveau de taux de ce document ne doit être lu comme une cotation réelle.
\end{abstract}

\newpage
{\setlength{\parskip}{0pt}\tableofcontents}
\newpage

% =============================================================================
\part{Théorie : la valorisation d'un swap OIS}
% =============================================================================

\section{Introduction}

Un \emph{swap OIS} (Overnight Indexed Swap) échange un taux fixe contre le taux overnight capitalisé
jour après jour sur la période. Depuis la crise de 2008, c'est l'instrument de référence pour
construire la courbe d'actualisation des produits collatéralisés : le taux overnight est le taux de
rémunération du collatéral, de sorte que la courbe OIS est le taux « sans risque » du marché
interbancaire~\cite{piterbarg2010,henrard2014}. En zone euro, l'indice est l'\estr{} (\emph{euro
short-term rate}), publié par la BCE depuis octobre~2019 à partir des opérations de dépôt non
garanties des banques ; il a remplacé l'EONIA, calculé d'octobre~2019 à janvier~2022 comme l'\estr{}
majoré de $8{,}5$~points de base.

\paragraph{Plan.} La partie~I est autonome : elle pose les notations (section~\ref{sec:notations}),
la capitalisation de l'\estr{} (section~\ref{sec:estr}), le swap (section~\ref{sec:swap}), sa
valorisation (section~\ref{sec:pricing}), la construction de la courbe (section~\ref{sec:bootstrap})
et le calcul des dates (section~\ref{sec:dates}). La partie~II décrit comment ces objets sont codés,
pourquoi les choix ont été faits ainsi, et comment le résultat est validé.

\paragraph{Périmètre.} La phase~1 couvre le swap OIS en euro, avec une courbe unique (actualisation
et projection sur la même courbe \estr{}). Les swaps contre Euribor, le multi-courbes et le calcul
de risque (DV01, key-rate, convexité) sont prévus en phase~2 (section~\ref{sec:limites}).

% -----------------------------------------------------------------------------
\section{Cadre et notations}\label{sec:notations}

\subsection{Facteurs d'actualisation, taux zéro et taux forward}

Soit $t_0$ la date de valorisation. Le \emph{facteur d'actualisation} (\emph{discount factor}) $\DF{t}$
est la valeur en $t_0$ d'un euro reçu en $t\ge t_0$, avec $\DF{t_0}=1$. La courbe est décrite par
ses valeurs en des dates \emph{piliers} et par un schéma d'interpolation. Le temps y est mesuré avec
la convention Act/365F : $\theta(d)=(d-t_0)/365$ en années.

\begin{definition}[Taux zéro, forward simple, forward instantané]
Pour une date $d>t_0$ et deux dates $s<e$ de longueur $\tau(s,e)$ selon une convention de décompte :
\begin{align}
  z(d) &= -\frac{\ln \DF{d}}{\theta(d)}\quad\text{(taux zéro, capitalisation continue)},\\
  F(s,e) &= \frac{1}{\tau(s,e)}\left(\frac{\DF{s}}{\DF{e}}-1\right)\quad\text{(forward simple)},
  \label{eq:forward}\\
  f(d) &= -\frac{\partial \ln \DF{d}}{\partial \theta}\quad\text{(forward instantané)}.
\end{align}
\end{definition}

Le forward simple $F(s,e)$ est le taux qui, appliqué à l'euro placé en $s$ jusqu'en $e$ sans risque,
rend l'opération équivalente à l'actualisation par la courbe : $1+F(s,e)\,\tau(s,e)=\DF{s}/\DF{e}$.

\subsection{Conventions de décompte}

La fraction d'année $\tau(d_1,d_2)$ dépend d'une convention (définitions ISDA~\cite{isda2006}).
Celles retenues sont :

\begin{center}
\begin{tabularx}{\linewidth}{@{}l>{\raggedright\arraybackslash}X@{}}
\toprule
Convention & Définition de $\tau(d_1,d_2)$\\
\midrule
Act/360 & $(d_2-d_1)/360$ (nombre exact de jours). Taux monétaires, \estr{}, jambes d'un OIS EUR.\\
Act/365F & $(d_2-d_1)/365$. Mesure du temps des courbes (interpolation).\\
30/360 (Bond Basis) &
  $\dfrac{360(Y_2-Y_1)+30(M_2-M_1)+(D_2'-D_1')}{360}$ avec $D_1'=\min(D_1,30)$ et $D_2'=30$ si
  $D_2=31$ et $D_1'=30$, sinon $D_2'=D_2$.\\
30E/360 (Eurobond) & Même formule avec $D_i'=\min(D_i,30)$.\\
\bottomrule
\end{tabularx}
\end{center}

\subsection{Calendriers et jours ouvrés}

Le calendrier TARGET compte comme jours fériés, en plus des week-ends : le 1\ier{} janvier, le
vendredi saint, le lundi de Pâques, le 1\ier{} mai, le 25 et le 26 décembre (les trois derniers
depuis 2000 ; le 31 décembre était aussi fermé en 1998, 1999 et 2001). Le dimanche de Pâques est
calculé par l'algorithme de Meeus--Jones--Butcher ; le vendredi saint est ce dimanche moins deux jours
et le lundi de Pâques plus un jour.

Une date tombant un jour non ouvré est déplacée selon une \emph{convention de jour ouvré}~:
\emph{Following} (jour ouvré suivant), \emph{Preceding} (précédent), ou \emph{Modified Following} :
jour ouvré suivant, sauf s'il change de mois, auquel cas on recule au jour ouvré précédent. On note
$\adj_c(d)$ la date ajustée selon la convention $c$ et $\adv(d,n)$ la date située $n$ jours ouvrés
après $d$.

% -----------------------------------------------------------------------------
\section{L'indice \estr{} et sa capitalisation}\label{sec:estr}

L'\estr{} est un taux d'intérêt d'un jour ouvré à l'autre. Le \emph{fixing} du jour ouvré $d$ est le
taux $r$ applicable de $d$ au jour ouvré suivant ; il est publié le jour ouvré suivant, donc connu
seulement a posteriori.

Sur une période de calcul $[S,E]$, soit $S=d_0<d_1<\dots<d_n=E$ les jours ouvrés TARGET, $r_i$ le
fixing du jour $d_i$ et $\delta_i=\tau_{360}(d_i,d_{i+1})$. Un investisseur qui replace chaque jour
son capital au taux overnight obtient le facteur de capitalisation
\begin{equation}
  C \;=\; \prod_{i=0}^{n-1}\bigl(1+r_i\,\delta_i\bigr),
  \qquad
  R \;=\; \frac{C-1}{\tau(S,E)},
  \label{eq:compounded}
\end{equation}
où $R$ est le \emph{taux overnight composé} de la période. Le montant d'intérêt sur un notionnel $N$
est $N(C-1)=N\,R\,\tau(S,E)$. Cette capitalisation est dite \emph{en arrears} : le taux n'est connu
qu'à la fin de la période. Dans ce document, on n'applique ni \emph{lookback}, ni \emph{lockout}, ni
\emph{observation shift} : chaque fixing est utilisé le jour même.

% -----------------------------------------------------------------------------
\section{Le swap OIS}\label{sec:swap}

\subsection{Structure}

Un swap OIS de notionnel $N$, de date de départ $S$ et de maturité $E$ échange deux jambes :
\begin{itemize}
  \item la \textbf{jambe fixe}, qui verse à chaque date $T^{\mathrm{fix}}_j$ le coupon
        $N\,K\,\tau^{\mathrm{fix}}_j$, où $K$ est le taux fixe ;
  \item la \textbf{jambe overnight}, qui verse à chaque date $T_j$ le coupon
        $N\,(C_j-1)$, où $C_j$ est le facteur \eqref{eq:compounded} de la période $[S_j,E_j]$.
\end{itemize}
Il n'y a pas d'échange de notionnel. Le détenteur d'un swap \emph{payeur} paie le taux fixe et reçoit
la jambe overnight ; un swap \emph{receveur} fait l'inverse. Dans tout le document, la valeur
actuelle (PV) est celle du point de vue du détenteur, et un montant reçu est compté positivement.

Chaque jambe est découpée en périodes $[S_j,E_j]$ (\emph{accrual}) ; le paiement a lieu à
$T_j=E_j$ augmenté d'un décalage de paiement (\emph{payment lag}) en jours ouvrés.

\subsection{Conventions EUR retenues}\label{sec:conv-eur}

Ces conventions ont été fixées avant l'implémentation (décision ADR-019) :

\begin{center}
\begin{tabularx}{\linewidth}{@{}>{\raggedright\arraybackslash}p{5.7cm}X@{}}
\toprule
Élément & Convention\\
\midrule
Calendrier & TARGET\\
Décalage de départ (\emph{spot lag}) & 2 jours ouvrés après la date de valorisation\\
Fréquence des deux jambes & Annuelle (une seule période si la maturité est d'un an ou moins)\\
Décompte des deux jambes & Act/360\\
Convention de jour ouvré & Modified Following (dates intermédiaires et maturité)\\
Règle de fin de mois & Appliquée si le départ est le dernier jour ouvré du mois\\
Génération de l'échéancier & \emph{Backward} depuis la maturité (période courte, ou \emph{stub}, en tête)\\
Décalage de paiement & 1 jour ouvré TARGET après la fin de la période\\
Capitalisation overnight & Quotidienne, en arrears, sans lookback, lockout ni observation shift\\
Dépôt overnight (calibration) & Départ en $t_0$, fin au jour ouvré suivant, Act/360, taux simple\\
\bottomrule
\end{tabularx}
\end{center}

% -----------------------------------------------------------------------------
\section{Valorisation d'un swap OIS}\label{sec:pricing}

\subsection{Hypothèse de courbe unique}

Le cadre est \emph{mono-courbe} : une seule courbe $P$ sert à actualiser les flux et à projeter les
fixings futurs. Cette hypothèse est cohérente pour un OIS : l'indice est le taux de la courbe
d'actualisation elle-même, de sorte que la capitalisation \eqref{eq:compounded} se réplique par un
dépôt overnight reconduit chaque jour. On pose donc, pour tout fixing futur,
$r_i=F(d_i,d_{i+1})$ au sens de \eqref{eq:forward} avec $\tau=\tau_{360}$.

\subsection{Télescopage du coupon overnight}

\begin{proposition}[Facteur de capitalisation projeté]\label{prop:telescope}
Si tous les fixings de $[S,E]$ sont projetés sur la courbe, alors
\[
  C \;=\; \prod_{i=0}^{n-1}\bigl(1+F(d_i,d_{i+1})\,\delta_i\bigr) \;=\; \frac{\DF{S}}{\DF{E}}.
\]
\end{proposition}

\begin{proof}
Par définition \eqref{eq:forward}, $1+F(d_i,d_{i+1})\,\delta_i=\DF{d_i}/\DF{d_{i+1}}$. Le produit est
télescopique : $\prod_{i=0}^{n-1}\DF{d_i}/\DF{d_{i+1}}=\DF{d_0}/\DF{d_n}$.
\end{proof}

La proposition~\ref{prop:telescope} est une identité exacte et non une approximation : elle dispense
de boucler sur les jours de la période (environ 250 par coupon annuel) et ramène le calcul d'un
coupon à deux évaluations de la courbe. Une vérification numérique est donnée à la
section~\ref{sec:exemple}.

\begin{remarque}
Si le paiement avait lieu en $E$ sans décalage, la valeur de la jambe overnight serait
$N\bigl(\DF{S}/\DF{E}-1\bigr)\DF{E}=N\bigl(\DF{S}-\DF{E}\bigr)$ : on retrouve la valeur d'un
\emph{floating rate note}. Le décalage de paiement $T_j>E_j$ rompt ce télescopage exact (voir
ci-dessous) et le montant, connu en $E_j$, est simplement actualisé de $T_j$.
\end{remarque}

\subsection{Valeur actuelle des deux jambes}

\begin{proposition}[PV des jambes]\label{prop:legs}
Soit $m$ coupons de dates d'accrual $[S_j,E_j]$ et de paiement $T_j$. Les valeurs actuelles,
non signées, des deux jambes sont
\begin{align}
  V_{\mathrm{fix}} &= N\,K\sum_{j}\tau^{\mathrm{fix}}_j\,\DF{T^{\mathrm{fix}}_j} \;=\; K\,A,
  \label{eq:vfix}\\
  V_{\mathrm{ON}} &= N\sum_{j=1}^{m}\left(\frac{\DF{S_j}}{\DF{E_j}}-1\right)\DF{T_j},
  \label{eq:von}
\end{align}
où $A=N\sum_j\tau^{\mathrm{fix}}_j\,\DF{T^{\mathrm{fix}}_j}$ est l'\emph{annuité}, exprimée en
unités de monnaie (le notionnel est inclus).
\end{proposition}

Lorsque $T_j=E_j$ et que les périodes sont contiguës ($S_{j+1}=E_j$), la somme \eqref{eq:von} se
télescope en $N\bigl(\DF{S_1}-\DF{E_m}\bigr)$. Avec un décalage de paiement, chaque terme est
multiplié par le facteur $\DF{T_j}/\DF{E_j}$, proche de $1$ ; le résultat n'a plus de forme fermée
aussi simple mais reste calculé exactement.

\subsection{Taux de par, annuité et PVBP}

\begin{definition}[Taux de par]
Le \emph{taux de par} (ou taux de swap) est le taux fixe $S_0$ pour lequel le swap a une valeur
nulle : $V_{\mathrm{fix}}=V_{\mathrm{ON}}$, soit
\begin{equation}
  S_0 \;=\; \frac{V_{\mathrm{ON}}}{A}.
  \label{eq:par}
\end{equation}
\end{definition}

La valeur d'un swap payeur de taux fixe $K$ est alors
\begin{equation}
  V_{\mathrm{payeur}} \;=\; V_{\mathrm{ON}}-V_{\mathrm{fix}} \;=\; A\,\bigl(S_0-K\bigr),
  \qquad V_{\mathrm{receveur}}=-V_{\mathrm{payeur}}.
  \label{eq:pv}
\end{equation}
Le taux de par est ce que le marché cote pour chaque maturité ; la formule \eqref{eq:pv} montre que
la valeur du swap est l'écart au par, multiplié par l'annuité. Le \emph{PVBP} (\emph{present value of
a basis point}) est la variation de la jambe fixe pour un point de base de taux fixe :
$\mathrm{PVBP}=A\times10^{-4}$, de sorte que $\partial V_{\mathrm{payeur}}/\partial K=-A$.

\subsection{Exemple chiffré : un swap payeur à 5 ans}\label{sec:exemple}

Prenons, à la date de valorisation \ValuationDate{}, un swap payeur de notionnel
$N=10\,000\,000\eur$, de taux fixe $K=\pct{2{,}25}$ et de maturité 5~ans, sur la courbe
\estr{} du snapshot synthétique. Le swap démarre au \FiveYStart{} (spot) et arrive à échéance le
\FiveYEnd{} ; il compte \FiveYCashflowCount{} flux futurs, soit 5 coupons par jambe. Le pricer
renvoie :

\begin{center}
\begin{tabular}{@{}lr@{}}
\toprule
Grandeur & Valeur\\
\midrule
Taux de par $S_0$ & \FiveYParRate\,\%\\
Annuité $A$ & \FiveYAnnuity\eur\\
PVBP $=A\times10^{-4}$ & \FiveYPVBP\eur\\
PV de la jambe fixe (payée) & \FiveYFixedLeg\eur\\
PV de la jambe overnight (reçue) & \FiveYOvernightLeg\eur\\
\textbf{PV du swap payeur} & \textbf{\FiveYPV\eur}\\
\bottomrule
\end{tabular}
\end{center}

La relation \eqref{eq:pv} se vérifie : $A\,(S_0-K)=\FiveYPVCheck\eur$, identique à la PV renvoyée.

\paragraph{Vérification du télescopage.} Pour le premier coupon overnight de ce swap
(du \TelStart{} au \TelEnd{}), le produit des \TelDays{} taux overnight implicites de la courbe,
calculés un à un sur les jours ouvrés TARGET, vaut $\TelProduct$, et le rapport
$\DF{S}/\DF{E}$ vaut $\TelRatio$ ; l'écart est de \TelDiff{}, c'est-à-dire de l'ordre de l'erreur
d'arrondi sur 257 multiplications en double précision.

\subsection{Swap en cours de vie}

Pour un swap déjà commencé à la date $t_0$, un coupon overnight $[S_j,E_j]$ peut être en partie
réalisé. Soit $d_k$ le premier jour d'observation dont le fixing n'est pas encore publié. Alors
\begin{equation}
  C_j \;=\; \underbrace{\prod_{d_i<d_k}\bigl(1+r_i\,\delta_i\bigr)}_{\text{fixings publiés}}
  \;\times\;\underbrace{\frac{\DF{d_k}}{\DF{E_j}}}_{\text{partie projetée}},
  \label{eq:seasoned}
\end{equation}
où la partie projetée utilise de nouveau la proposition~\ref{prop:telescope}. Les règles suivantes
s'appliquent (décision ADR-020) :
\begin{itemize}
  \item les fixings \emph{strictement antérieurs} à $t_0$ sont obligatoires ; leur absence est une
        erreur de données de marché ;
  \item le fixing du jour $t_0$ est utilisé s'il est fourni, sinon il est projeté ;
  \item un flux payé en $t_0$ ou avant est échu et n'entre pas dans la PV.
\end{itemize}

% -----------------------------------------------------------------------------
\section{Construction de la courbe par bootstrap}\label{sec:bootstrap}

\subsection{Principe}

Le marché cote les taux de par de swaps OIS de maturités $1\text{W},2\text{W},\dots,30\text{Y}$ ainsi
que le taux du dépôt overnight. La courbe $P$ est déterminée par ses valeurs aux \emph{piliers}
$T_1<\dots<T_K$ : le pilier d'un dépôt est sa maturité, celui d'un OIS est sa \emph{dernière date de
paiement}, qui est la date la plus lointaine dont dépend son taux de par.

\begin{definition}[Équations de calibration]
Pour chaque instrument $k$ de cotation $q_k$, on cherche $P_k=\DF{T_k}$ tel que la cotation
implicite de la courbe égale la cotation de marché :
\begin{itemize}
  \item dépôt overnight, de $s=t_0$ à $e$ : $q=\dfrac{1}{\tau_{360}(s,e)}\Bigl(\dfrac{1}{\DF{e}}-1\Bigr)$,
        d'où la solution fermée $\DF{e}=1/\bigl(1+q\,\tau_{360}(s,e)\bigr)$ ;
  \item OIS de maturité $k$ : $S_0\bigl(P_1,\dots,P_k\bigr)=q_k$, avec $S_0$ donné par
        \eqref{eq:par}.
\end{itemize}
\end{definition}

L'équation de l'OIS n'a pas de solution fermée : les facteurs d'actualisation aux dates
intermédiaires sont interpolés entre piliers, et le taux de par dépend de $P_k$ de manière non
linéaire. On la résout donc numériquement, pilier par pilier, dans l'ordre croissant des maturités.

\paragraph{Pourquoi une seule passe suffit.} Avec une interpolation \emph{locale} (la valeur
interpolée sur un segment $[T_{k-1},T_k]$ ne dépend que des nœuds $P_{k-1}$ et $P_k$), le taux de par
de l'instrument $k$ ne dépend que des nœuds $P_1,\dots,P_k$, puisque tous ses flux sont datés au plus
tard en $T_k$. Les nœuds déjà trouvés ne sont donc jamais remis en cause : l'algorithme est séquentiel
et exact, sans itération globale.

\subsection{Interpolation}

Soit $\theta_k=\theta(T_k)$, avec $\theta_0=0$. Deux schémas sont disponibles :
\begin{itemize}
  \item \textbf{Log-linéaire sur les facteurs d'actualisation} (par défaut) : pour
        $\theta\in[\theta_{k-1},\theta_k]$,
        \[\ln P=\ln P_{k-1}+\frac{\theta-\theta_{k-1}}{\theta_k-\theta_{k-1}}\bigl(\ln P_k-\ln P_{k-1}\bigr),\]
        ce qui revient à un forward instantané $f$ \emph{constant par morceaux}. C'est le standard des
        courbes OIS~\cite{hagan2006}.
  \item \textbf{Linéaire sur le taux zéro} : $z$ est linéaire en $\theta$ entre les nœuds, avec
        $z(\theta\!=\!0)$ pris égal au zéro du premier pilier.
\end{itemize}
Au-delà du dernier pilier $\theta_K$, les deux schémas prolongent la courbe à \emph{forward instantané
constant}. Pour le schéma log-linéaire, c'est le forward du dernier segment ; pour le schéma
linéaire sur le zéro, c'est $f_K=z_K+\theta_K\,z'(\theta_K)$, d'où
$-\ln P=z_K\,\theta_K+f_K\,(\theta-\theta_K)$.

\subsection{Résolution numérique}

Pour le pilier $k$, on cherche la racine de
$g(P)=S_0(P_1,\dots,P_{k-1},P)-q_k$ dans l'intervalle
\[
  \Bigl[\,P_{k-1}\,e^{-1{,}0\,\Delta\theta},\;P_{k-1}\,e^{+0{,}5\,\Delta\theta}\,\Bigr],
  \qquad \Delta\theta=\theta_k-\theta_{k-1},
\]
qui correspond à des forwards continus compris entre $-50\,\%$ et $+100\,\%$ : large pour tout
marché réaliste, mais suffisamment serré pour que $g$ change de signe. Plus $P$ est faible, plus les
taux implicites sont élevés ; $g$ est en pratique strictement décroissante, ce qui rend la racine
unique dans l'intervalle.

\paragraph{Méthode de Brent.} Brent~\cite{brent1973} combine la bissection (robuste, convergence
linéaire), la sécante et l'interpolation quadratique inverse (convergence superlinéaire). À chaque
itération, le pas d'interpolation n'est accepté que s'il reste à l'intérieur de l'encadrement et
réduit suffisamment l'intervalle ; sinon on bissecte. La convergence est garantie dès que
$g(P_{\min})$ et $g(P_{\max})$ sont de signes opposés.

\paragraph{Méthode de Newton protégée.} Une variante de Newton--Raphson, avec dérivée par différences
finies centrées, est disponible en option ; tout pas sortant de l'encadrement est remplacé par une
bissection, ce qui lui conserve la robustesse de la bissection.

\paragraph{Critère d'échec.} Un pilier dont le résidu final $\lvert S_0-q_k\rvert$ dépasse
$10^{-12}$ (en taux décimal) est une erreur de calibration, qui identifie le pilier fautif et son
résidu.

% -----------------------------------------------------------------------------
\section{Dates : spot, maturité, échéancier}\label{sec:dates}

\paragraph{Date spot.} Avec le calendrier $\mathcal{C}$ et le décalage $\ell=2$ jours ouvrés :
$t_{\mathrm{spot}}=\adv\bigl(\adj_{F}(t_0),\,\ell\bigr)$, où $F$ désigne \emph{Following}.

\paragraph{Maturité.} Si la règle de fin de mois s'applique (le spot est le dernier jour ouvré de
son mois), la maturité d'un swap de tenor $n$ mois est le dernier jour ouvré du mois situé $n$ mois
plus tard. Sinon, c'est $t_{\mathrm{spot}}+n$ mois, ajustée ensuite par l'échéancier.

\paragraph{Échéancier (génération \emph{backward}).}
\begin{enumerate}
  \item À partir de la date de fin, on recule par pas de la fréquence (1~an) tant que la date obtenue
        n'est pas antérieure à la date de départ ; le reliquat plus court que la fréquence est une période courte
        (\emph{stub}), placée en tête.
  \item Si la règle de fin de mois s'applique, les dates intermédiaires et la date de fin sont ramenées
        au dernier jour ouvré de leur mois ; la date de départ est seulement ajustée.
  \item Sinon, chaque date est ajustée en Modified Following.
  \item La date de paiement d'une période est la date de fin ajustée, avancée de 1 jour ouvré
        TARGET.
\end{enumerate}

\paragraph{Illustration.} Pour un swap payeur à 2~ans de notionnel $10\,000\,000\eur$ et de taux fixe
$\pct{2{,}00}$, le tableau~\ref{tab:cashflows} donne l'échéancier et les flux calculés. La
deuxième période se termine le 2028-10-23 et non le 2028-10-21, qui est un samedi (Modified
Following), et le paiement a lieu le jour ouvré suivant. Le taux de la jambe overnight est le taux
composé \eqref{eq:compounded} de chaque période ; le taux de par vaut \TwoYParRate\,\%, l'annuité
\TwoYAnnuity\eur{} et la PV du swap payeur \TwoYPV\eur.

\begin{table}[ht]
\centering
\caption{Flux d'un swap payeur à 2 ans (notionnel $10^7\eur$, $K=\pct{2{,}00}$). Montants signés :
positifs si reçus.}\label{tab:cashflows}
\setlength{\tabcolsep}{3.2pt}
\resizebox{\linewidth}{!}{\begin{tabular}{@{}lccccrrr@{}}
\toprule
Jambe & Début & Fin & Paiement & Taux (\%) & Montant (\texteuro) & DF & PV (\texteuro) \\
\midrule
Fixe (pay) & 2026-10-21 & 2027-10-21 & 2027-10-22 & 2{,}0000 & $-$202\,777{,}78 & 0{,}981090 & $-$198\,943{,}29 \\
Fixe (pay) & 2027-10-21 & 2028-10-23 & 2028-10-24 & 2{,}0000 & $-$204\,444{,}44 & 0{,}960462 & $-$196\,361{,}07 \\
Overnight (receive) & 2026-10-21 & 2027-10-21 & 2027-10-22 & 1{,}8850 & 191\,118{,}06 & 0{,}981090 & 187\,504{,}05 \\
Overnight (receive) & 2027-10-21 & 2028-10-23 & 2028-10-24 & 2{,}1004 & 214\,708{,}28 & 0{,}960462 & 206\,219{,}09 \\
\bottomrule
\end{tabular}
}
\end{table}

% =============================================================================
\part{Implémentation et choix techniques}
% =============================================================================

\section{Vue d'ensemble et principes}\label{sec:archi}

\subsection{Architecture en couches}

Le projet est un paquet Python unique, \code{rate_derivatives_pricer}, organisé en cinq couches dont
les dépendances vont dans \emph{un seul sens} (figure~\ref{fig:layers}).

\begin{figure}[ht]
\centering
\begin{tikzpicture}[node distance=0.5cm]
  \node[layer] (ui) {\textbf{ui}\\[-1pt]{\small Streamlit + Plotly}};
  \node[layer, below=of ui] (sdk) {\textbf{sdk}\\[-1pt]{\small client typé \texttt{httpx}}};
  \node[layer, below=of sdk] (api) {\textbf{api}\\[-1pt]{\small FastAPI, schémas Pydantic}};
  \node[layer, below=of api] (srv) {\textbf{services}\\[-1pt]{\small orchestration, cache, mappers}};
  \node[layer, below=of srv, fill=green!10] (core) {\textbf{core}\\[-1pt]{\small bibliothèque quantitative (NumPy)}};
  \draw[flow] (ui) -- (sdk);
  \draw[flow] (sdk) -- (api) node[midway, right=2pt, font=\scriptsize] {HTTP / JSON};
  \draw[flow] (api) -- (srv);
  \draw[flow] (srv) -- (core);
  \node[right=2.2cm of srv, align=left, font=\small, text width=5.2cm] (rule)
    {\textbf{Règle imposée par \texttt{import-linter}} : \texttt{core} n'importe rien des
     autres couches, ni aucune bibliothèque web, SciPy ou QuantLib.};
  \draw[dashed, black!50] (rule.west) -- ++(-0.9,0) |- (core.east);
\end{tikzpicture}
\caption{Couches du projet. Les flèches indiquent le sens de la dépendance.}\label{fig:layers}
\end{figure}

La règle est vérifiée en intégration continue par cinq contrats \code{import-linter} :
(1)~la hiérarchie $\text{ui}\to\text{sdk}\to\text{api}\to\text{services}\to\text{core}$ ;
(2)~\code{ui} n'atteint jamais \code{services} ni \code{core}, même indirectement ;
(3)~\code{ui} n'importe jamais \code{api} directement ; (4)~\code{sdk} ne descend pas sous l'API ;
(5)~\code{core} ne dépend d'aucun paquet web, d'interface, de SciPy ou de QuantLib.
Un import interdit fait échouer la CI.

\subsection{Principes transversaux}

\begin{itemize}
  \item \textbf{Immutabilité.} Instruments, cotations, échéanciers, snapshots et résultats sont des
        \code{dataclass} figées (\code{frozen=True}, \code{slots=True}). Un choc de risque produira de nouveaux objets
        plutôt que de modifier les anciens.
  \item \textbf{Pas d'état global.} La date de valorisation est une donnée explicite du
        \code{MarketSnapshot}. Le \code{Settings.evaluationDate} global de QuantLib n'est volontairement
        pas reproduit (ADR-009) ; les tests alignent QuantLib sur la date du snapshot.
  \item \textbf{Données et comportement séparés.} Un instrument décrit un contrat, un moteur le
        valorise, un schéma d'API le transporte. Les instruments ne portent pas leur moteur,
        contrairement à QuantLib (ADR-007).
  \item \textbf{Conventions en configuration.} Aucune convention de marché n'est codée en dur
        dans le calcul : elles sont lues dans un fichier JSON par devise (ADR-002, ADR-019).
  \item \textbf{Erreurs typées.} Une exception par mode d'échec, traduite en code HTTP stable.
  \item \textbf{Le cœur émet les logs, l'API les configure.} \code{core} n'utilise que le module
        \code{logging} standard.
\end{itemize}

\subsection{Pipeline de calcul}

Le flux est fonctionnel : le snapshot est transformé en environnement de marché par le bootstrap, puis l'instrument est valorisé dans cet environnement par un moteur sans état (figure~\ref{fig:pipeline}). Le risque se ramènera à « modifier le snapshot, puis rejouer le pipeline ».

\begin{figure}[ht]
\centering
\begin{tikzpicture}[node distance=0.55cm]
  \node[data] (snap) {\code{MarketSnapshot}\\{\scriptsize date, cotations, fixings}};
  \node[proc, right=of snap] (builder) {\code{CurveBuilder}\\{\scriptsize bootstrap}};
  \node[data, right=of builder] (env) {\code{MarketEnvironment}\\{\scriptsize courbe, fixings}};
  \node[proc, below=1.1cm of env] (engine) {\code{DiscountingSwapEngine}};
  \node[data, left=of engine] (instr) {\code{OisSwap}\\{\scriptsize jambes, coupons}};
  \node[data, below=0.9cm of engine] (res) {\code{PricingResult}\\{\scriptsize PV, par, annuité}};
  \draw[flow] (snap) -- (builder);
  \draw[flow] (builder) -- (env);
  \draw[flow] (env) -- (engine);
  \draw[flow] (instr) -- (engine);
  \draw[flow] (engine) -- (res);
\end{tikzpicture}
\caption{Pipeline fonctionnel du cœur : chaque étape prend des entrées et renvoie une sortie.}
\label{fig:pipeline}
\end{figure}

\section{Le cœur quantitatif}\label{sec:core}

Le cœur ne dépend que de NumPy et de la bibliothèque standard. Il compte les sous-paquets
suivants (la liste complète des modules figure en annexe~\ref{app:modules}).

\subsection{Conventions : calendriers, décomptes, tenors}

\code{core/conventions} implémente les objets de la section~\ref{sec:notations}.
\code{Calendar} est une classe abstraite dont dérivent \code{TargetCalendar},
\code{NullCalendar}, \code{WeekendsOnlyCalendar} et \code{JointCalendar} (union des jours fériés).
Les jours fériés TARGET sont calculés par règles (Pâques par l'algorithme de Meeus) et mis en cache,
car le calendrier est interrogé jour par jour. \code{Tenor} (« 3M », « 10Y », « ON ») et
\code{Frequency} sont des objets valeur ; l'addition de mois ramène le jour à la longueur du mois
cible (31~janvier + 1~mois = 28 ou 29~février).

Les classes \code{Act360}, \code{Act365Fixed}, \code{Thirty360} et \code{Thirty360European} se
déclarent dans un registre par un décorateur :

\begin{lstlisting}[caption={Enregistrement d'un composant par décorateur (principe Open/Closed).}]
@DAY_COUNTS.register("ACT/360")
@dataclass(frozen=True, slots=True)
class Act360(DayCountConvention):
    name: ClassVar[str] = "ACT/360"

    def day_count(self, start: date, end: date) -> int:
        return (end - start).days

    def year_fraction(self, start: date, end: date) -> float:
        return self.day_count(start, end) / DAYS_PER_YEAR_360
\end{lstlisting}

Ajouter un composant revient donc à ajouter un fichier, sans modifier de ligne existante.
Deux variantes 30/360 explicites (Bond Basis et Eurobond) remplacent une classe \code{Thirty360}
ambiguë (ADR-019).

\paragraph{Conventions de devise.} \code{CurrencyConventions} regroupe calendrier, \emph{spot lag},
mesure du temps des courbes, indice overnight, conventions du dépôt et de l'OIS. Elle est construite
à partir d'un fichier JSON par devise (listing~\ref{lst:eur}) ; l'ajout de l'USD en phase~3 se fera par
un nouveau fichier, sans toucher au code de calcul.

\begin{lstlisting}[language=json,caption={\code{core/conventions/data/EUR.json}.},label={lst:eur}]
{
  "currency": "EUR",
  "calendar": "TARGET",
  "spot_lag_days": 2,
  "curve_day_count": "ACT/365F",
  "overnight_index": { "name": "ESTR", "calendar": "TARGET", "day_count": "ACT/360" },
  "deposit": { "settlement_days": 0, "business_day_convention": "MODIFIED_FOLLOWING",
               "end_of_month": false, "day_count": "ACT/360" },
  "ois": { "fixed_frequency": "ANNUAL", "fixed_day_count": "ACT/360",
           "floating_frequency": "ANNUAL", "floating_day_count": "ACT/360",
           "business_day_convention": "MODIFIED_FOLLOWING", "end_of_month": true,
           "rule": "BACKWARD", "payment_lag_days": 1 }
}
\end{lstlisting}

\subsection{Échéanciers}

\code{ScheduleGenerator} transforme (départ, fin, fréquence) en une suite de \code{SchedulePeriod}
(début et fin d'accrual, date de paiement, indicateur de période courte). Il suit pas à pas
l'algorithme de la section~\ref{sec:dates}. La comparaison date par date avec QuantLib a révélé un
détail de convention : en génération \emph{backward} avec la règle de fin de mois, QuantLib ne
déplace \emph{pas} la date de départ en fin de mois, seule la date « pivot » l'est. Notre première version
déplaçait les deux ; le test de validation l'a détecté et le générateur a été corrigé.

\subsection{Données de marché}

\begin{itemize}
  \item \code{Quote} (\code{DepositQuote}, \code{OisSwapQuote}) porte \emph{toujours une unité}
        explicite (\code{percent}, \code{decimal} ou \code{bp}) ; la propriété \code{rate} convertit en
        décimal. Confondre $2{,}5$ et $0{,}025$ est un bug de production classique, que l'unité
        obligatoire rend impossible.
  \item \code{MarketSnapshot} regroupe la date de valorisation, les cotations et les fixings.
        Il refuse les cotations en double et les fixings postérieurs à la date de valorisation.
        C'est l'\emph{unité de choc} du futur calcul de risque.
  \item \code{FileMarketDataProvider} lit \code{data/market/<nom>.json}. Les noms sont filtrés par
        expression régulière pour interdire toute remontée d'arborescence (\code{../}). Le format
        canonique du snapshot sert aussi de base à l'empreinte du cache.
\end{itemize}

\subsection{Courbes et interpolation}

\code{DiscountCurve} est définie par des nœuds (date, facteur d'actualisation), un interpolateur et
une convention de décompte interne. Le temps de chaque nœud est converti une fois, à la
construction, en tableaux NumPy mis en cache dans l'objet figé. Les interpolateurs
(\code{LogLinearOnDiscountFactor}, \code{LinearOnZeroRate}) sont deux stratégies interchangeables
derrière une interface \code{Interpolator} qui travaille sur des tableaux : l'évaluation d'une courbe
sur une grille de dates est vectorisée.

\subsection{Solveurs}

\code{BrentSolver} et \code{NewtonSolver} implémentent l'interface \code{Solver} (pas de SciPy dans
le cœur, ADR-004). Ils renvoient un \code{SolverResult} (racine, résidu, itérations, convergence)
plutôt que de lever une exception : c'est le bootstrap qui décide de la gravité d'un échec. La
tolérance absolue par défaut est de $10^{-15}$ sur le facteur d'actualisation, et le nombre
maximal d'itérations de 100.

\subsection{Instruments de calibration et bootstrap}

L'abstraction \code{CalibrationInstrument} est l'équivalent du \code{RateHelper} de QuantLib : pour
une courbe candidate, elle calcule la cotation implicite et indique le pilier qu'elle détermine.
\code{DepositCalibration} renvoie le forward simple entre début et fin ; \code{OisCalibration}
construit un swap de notionnel~1 aux conventions de la devise et renvoie son taux de par. Cette
abstraction rend le bootstrap générique : ajouter un instrument (un FRA, par exemple) n'exige
aucune modification de \code{CurveBuilder}.

\code{CurveBuilder.build(snapshot)} trie les instruments par pilier, refuse deux instruments sur le
même pilier, puis résout chaque nœud avec le solveur. Le cœur de la boucle est le suivant
(extrait, journalisation omise) :

\begin{lstlisting}[caption={Calibration d'un pilier (\code{core/market/bootstrap.py}).}]
market_quote = instrument.quote.rate
previous_df = discount_factors[-1] if discount_factors else 1.0
dt = template.time(pillar) - previous_time

def objective(discount_factor: float) -> float:
    candidate = template.with_pillars(dates, (*discount_factors, discount_factor))
    return instrument.implied_quote(environment(candidate)) - market_quote

result = self._solver.solve(
    objective,
    lower=previous_df * math.exp(-BRACKET_MAX_FORWARD_RATE * dt),
    upper=previous_df * math.exp(-BRACKET_MIN_FORWARD_RATE * dt),
    initial_guess=previous_df * math.exp(-market_quote * dt),
)
if not result.converged or not abs(result.residual) <= CALIBRATION_RESIDUAL_TOLERANCE:
    raise CalibrationError(
        f"bootstrap failed on pillar {instrument.label}",
        pillar=instrument.label, residual=result.residual,
    )
\end{lstlisting}

Le résultat est un \code{MarketEnvironment} : la courbe (qui sert à la fois d'actualisation et de
projection de l'indice \code{ESTR}), les fixings et un \code{CalibrationReport} qui détaille chaque
pilier (cotation, cotation implicite, résidu, itérations). Sur le snapshot de référence
(\QuoteCount{} cotations), la calibration compte \PillarCount{} piliers, au plus \MaxIterations{}
itérations par pilier et un résidu maximal de \MaxResidual{}, en environ 0,1~s. La
figure~\ref{fig:curve} montre la courbe obtenue.

\begin{figure}[ht]
\centering
\ref{curvelegend}\\[4pt]
\begin{tikzpicture}
\begin{axis}[
  name=left, width=0.4\linewidth, height=6.2cm,
  xlabel={Maturité (années)}, ylabel={Taux (\%)}, xmin=0, xmax=30,
  legend to name=curvelegend, legend columns=3,
]
\addplot[blue!70!black, thick] table[x=t, y=zero]{data/curve.dat};
\addlegendentry{Zéro continu (Act/365F)}
\addplot[orange!85!black, thick, dashed] table[x=t, y=forward]{data/curve.dat};
\addlegendentry{Forward 1M (Act/360)}
\addplot[only marks, mark=*, mark size=1.4pt, black] table[x=t, y=zero]{data/pillars.dat};
\addlegendentry{Piliers calibrés}
\end{axis}
\begin{axis}[
  at={(left.east)}, anchor=west, xshift=1.3cm, width=0.4\linewidth, height=6.2cm,
  xlabel={Maturité (années)}, ylabel={Facteur d'actualisation},
  xmin=0, xmax=30, ymin=0.4, ymax=1,
]
\addplot[blue!70!black, thick] table[x=t, y=df]{data/curve.dat};
\end{axis}
\end{tikzpicture}
\caption{Courbe \estr{} bootstrapée sur le snapshot synthétique du \ValuationDate{}, tracée à partir des
données produites par le pricer.}\label{fig:curve}
\end{figure}

\subsection{Instruments, jambes et coupons}

\begin{itemize}
  \item \code{Cashflow} est une classe abstraite : chaque flux \emph{calcule lui-même son montant}
        à partir du \code{MarketEnvironment} (ADR-007). \code{Coupon} (notionnel, dates d'accrual,
        date de paiement, décompte) en dérive, avec \code{FixedCoupon} et
        \code{OvernightCompoundedCoupon}.
  \item \code{Leg} regroupe des flux, un notionnel et un sens (\code{PayReceive}, signe $\pm1$) :
        la composition plutôt que l'héritage.
  \item \code{OisSwap} décrit le contrat ; \code{SwapDirection} (\code{pay_fixed} ou
        \code{receive_fixed}) est nommé d'après la jambe fixe.
  \item \code{factories.py} construit un swap à partir des conventions : \code{spot_date},
        \code{ois_maturity_date} (qui reproduit la logique de \code{MakeOIS}, règle de fin de mois
        comprise), \code{build_ois_swap} (deux échéanciers puis deux jambes) et la factory
        \code{ois_swap_factory} enregistrée sous la clé \code{ois_swap}.
\end{itemize}

Le coupon overnight est la pièce la plus délicate. Il applique l'équation \eqref{eq:seasoned} :
boucle jour par jour sur la partie réalisée, télescopage sur la partie projetée
(proposition~\ref{prop:telescope}).

\begin{lstlisting}[caption={Facteur de capitalisation d'un coupon overnight (\code{cashflows.py}, docstrings omises).}]
def compound_factor(self, env: MarketEnvironment) -> float:
    valuation_date = env.valuation_date
    curve = env.projection_curve(self.index.name)
    if self.accrual_start > valuation_date:
        # Entirely in the future: telescoped projection.
        start_df, end_df = curve.discount_factors([self.accrual_start, self.accrual_end])
        return float(start_df / end_df)

    observation_dates = [d for d in self.index.fixing_calendar.business_days(
        self.accrual_start, self.accrual_end) if d < self.accrual_end]
    period_ends = [*observation_dates[1:], self.accrual_end]

    factor, i, n = 1.0, 0, len(observation_dates)
    # Realised part: fixings strictly before the valuation date are mandatory.
    while i < n and observation_dates[i] < valuation_date:
        factor *= self._realised_growth(env, observation_dates[i], period_ends[i])
        i += 1
    # Fixing of the valuation date itself: used if already known, projected otherwise.
    if i < n and observation_dates[i] == valuation_date:
        todays_fixing = env.fixing(self.index.name, valuation_date)
        if todays_fixing is not None:
            factor *= 1.0 + todays_fixing * self.index.day_count.year_fraction(
                observation_dates[i], period_ends[i])
            i += 1
    # Projected part, telescoped on the projection curve.
    if i < n:
        start_df, end_df = curve.discount_factors([observation_dates[i], self.accrual_end])
        factor *= float(start_df / end_df)
    return factor
\end{lstlisting}

Un fixing passé manquant lève une \code{MarketDataError} qui précise l'indice et la date, et se
traduit côté API par une erreur $422$ de code \code{MARKET_DATA_INVALID}.

\subsection{Moteur de valorisation}

\code{PricingEngine} est une interface ; \code{engine_for(type, method)} choisit l'implémentation
dans le registre \code{ENGINES}. Les moteurs sont sans état. \code{DiscountingSwapEngine} implémente
les formules \eqref{eq:vfix}--\eqref{eq:pv} :

\begin{lstlisting}[caption={Fonctions de base du moteur (\code{engines/discounting.py}).}]
def _alive_coupons(leg: Leg, env: MarketEnvironment) -> list[Coupon]:
    # Coupons of the leg paid strictly after the valuation date.
    return [cf for cf in leg.cashflows
            if isinstance(cf, Coupon) and cf.payment_date > env.valuation_date]

def par_rate(self, instrument: Instrument, env: MarketEnvironment) -> float:
    swap = self._check(instrument)
    annuity = _annuity(swap.fixed_leg, env)          # sum N * tau_i * P(pay_i)
    if annuity == 0.0:
        raise PricingError("par rate undefined: no remaining fixed coupon")
    return _unsigned_leg_value(swap.overnight_leg, env) / annuity
\end{lstlisting}

La méthode \code{price} renvoie un \code{PricingResult}, \emph{jamais un simple flottant} : PV
total, PV par jambe, tableau des flux (dates, fraction d'année, taux, montant signé, facteur
d'actualisation, PV), taux de par, annuité, PVBP et diagnostics (moteur, durée, version du cœur,
avertissements). Si aucun coupon fixe ne subsiste, le taux de par est indéfini : il vaut \code{NaN}
dans le cœur, un avertissement est émis, et l'API renvoie \code{null}. La méthode \code{par_rate}
est la version allégée, sans tableau de flux, réutilisée par le bootstrap.

\subsection{Registres et erreurs}

Six registres (\code{CALENDARS}, \code{DAY_COUNTS}, \code{INTERPOLATORS}, \code{INSTRUMENTS},
\code{ENGINES}, \code{CONVENTIONS}) associent une clé stable à un composant. Leurs modules sont
importés \emph{paresseusement} à la première consultation : un module de calendrier importe le
registre pour s'y enregistrer, donc le registre ne peut pas importer les calendriers au chargement
sans créer d'import circulaire. Le chargeur est commun à tous les registres et réentrant. Le
catalogue de l'API (\code{/v1/catalog}) est généré à partir de ces registres.

{\sloppy La hiérarchie d'erreurs \code{RatesPricerError} comporte une classe par mode d'échec
(\code{ConventionError}, \code{ScheduleError}, \code{MarketDataError}, \code{CalibrationError},
\code{UnsupportedInstrumentError}, \code{PricingError}, \code{RiskComputationError}). Chacune porte
un message et des détails structurés, que l'API met en correspondance avec un code HTTP
(tableau~\ref{tab:errors}).\par}

% -----------------------------------------------------------------------------
\section{Services, API, SDK et interface}\label{sec:services}

\subsection{Couche services}

\code{PricerService} est la façade appelée par l'API. Ses entrées sont des dataclasses figées
(\code{SnapshotReference} ou \code{InlineMarket} pour le marché, \code{TradeInput},
\code{CalculationOptions}) : Pydantic reste cantonné à \code{api} et \code{sdk} (ADR-021). Pour
chaque calcul, le service :
\begin{enumerate}
  \item charge le snapshot, par référence (fichier) ou par valeur (corps de la requête) ;
  \item calcule l'\emph{empreinte} du marché et consulte le cache de courbes ;
  \item en cas d'absence, bootstrape la courbe et la met en cache ;
  \item construit l'instrument avec la factory, le valorise avec le moteur ;
  \item renvoie le résultat accompagné de diagnostics (durées par étape, \emph{hit} ou \emph{miss} du
        cache, avertissements, empreintes) et émet un événement d'audit.
\end{enumerate}

\paragraph{Cache de courbes.} \code{LruCurveCache} est un cache en mémoire, \emph{least recently
used}, à capacité configurable (32 par défaut) et à durée de vie optionnelle, protégé par un verrou
(l'API sert plusieurs requêtes dans des threads). Il tient des compteurs (\emph{hits}, \emph{misses},
évictions) et accepte une horloge injectable, ce qui permet de tester l'expiration sans attente. Il
est \emph{par processus} : avec plusieurs \emph{workers}, chacun a son cache. Une interface
\code{CurveCache} permettra de le remplacer par Redis (ADR-005).

\paragraph{Clé du cache.} La clé est le SHA-256 de la sérialisation JSON canonique (clés triées)
des données du snapshot (date, devise, cotations avec leurs unités, fixings), de l'interpolation,
du solveur et de la version du cœur. Le \emph{nom} du snapshot en est exclu : les mêmes données
passées par référence ou par valeur partagent la même entrée (ADR-026). Une même empreinte d'entrée et
une même version doivent donner exactement le même résultat ; un PV différent révèle alors un bug.

\subsection{API HTTP}

L'API FastAPI expose, en phase~1, les routes du tableau~\ref{tab:routes}. Principes du contrat :

\begin{table}[ht]
\centering
\caption{Routes de l'API en phase~1.}\label{tab:routes}
\begin{tabularx}{\linewidth}{@{}>{\raggedright\arraybackslash}p{5.3cm}X@{}}
\toprule
Route & Rôle\\
\midrule
\code{GET /health} & Liveness : le processus répond.\\
\code{GET /ready} & Readiness : conventions et registres chargés (503 sinon).\\
\code{GET /v1/version} & Versions de l'API, du cœur, de Python et de NumPy ; commit déployé.\\
\code{GET /v1/catalog} & Types d'instruments, interpolations, devises, schémas des trades, générés
  depuis les registres.\\
\code{GET /v1/market-snapshots} & Liste des snapshots ; le contenu d'un snapshot est sous
  \code{/v1/market-snapshots/<nom>}.\\
\code{POST /v1/curves/build} & Bootstrap : nœuds calibrés, grille de la courbe, diagnostics de
  convergence.\\
\code{POST /v1/price} & Valorisation d'un trade.\\
\bottomrule
\end{tabularx}
\end{table}

\begin{itemize}
  \item \textbf{Sans état} : chaque requête porte ou référence ses données de marché et son trade.
  \item \textbf{Versionnée par l'URL} (\code{/v1}) : un changement incompatible crée \code{/v2}.
  \item \textbf{DTO distincts du domaine} : les schémas Pydantic (\code{api/schemas.py}) ne dépendent
        que de Pydantic ; les champs inconnus sont refusés.
  \item \textbf{Unités explicites} : dates ISO~8601, taux en décimal, cotations avec unité.
  \item \textbf{Enveloppe permanente} : chaque réponse contient \code{request_id}, \code{result} et un
        bloc \code{diagnostics}.
\end{itemize}

\begin{lstlisting}[language=json,caption={Requête \code{POST /v1/price} et réponse (abrégée).}]
POST /v1/price
{ "market":  { "snapshot": "EUR_2026-10-19" },
  "trade":   { "type": "ois_swap",
               "params": { "notional": 10000000, "fixed_rate": 0.0225,
                           "direction": "pay_fixed", "maturity": "5Y" } },
  "options": { "interpolation": "log_linear_df" } }

200 OK
{ "request_id": "smoke-1",
  "result": { "instrument_type": "ois_swap", "currency": "EUR",
              "present_value": 1902.5377349308692, "par_rate": 0.02254000000000001,
              "annuity": 47563443.373254724, "legs": [ ... ], "cashflows": [ ... ] },
  "diagnostics": { "curve_cache": "miss", "timings_ms": { "curve_build": 99.4, "pricing": 0.7 },
                   "versions": { "api": "0.1.0", "core": "0.1.0" }, "warnings": [],
                   "input_hashes": { "market": "8295394e...", "trade": "dc6b132c..." } } }
\end{lstlisting}

\paragraph{Gestion des erreurs.} Un format unique est utilisé partout,
\code{{"request_id", "error": {"code", "message", "details"}}}, avec un code machine stable
(tableau~\ref{tab:errors}). Une exception inattendue produit une réponse générique \code{INTERNAL_ERROR}
sans trace ; la trace complète est journalisée avec le \code{request_id}.

\begin{table}[ht]
\centering
\caption{Correspondance entre erreurs du cœur et réponses HTTP.}\label{tab:errors}
\begin{tabularx}{\linewidth}{@{}>{\raggedright\arraybackslash}X c l@{}}
\toprule
Source & HTTP & Code\\
\midrule
Validation Pydantic (JSON mal formé, champ manquant ou inconnu) & 422 & \code{VALIDATION_ERROR}\\
\code{MarketDataError} & 422 & \code{MARKET_DATA_INVALID}\\
\code{ConventionError}, \code{ScheduleError} & 422 & \code{INVALID_TRADE_DEFINITION}\\
\code{UnsupportedInstrumentError} & 400 & \code{UNSUPPORTED_INSTRUMENT}\\
\code{CalibrationError} (pilier et résidu) & 422 & \code{CALIBRATION_FAILED}\\
\code{MarketDataNotFoundError}, route inconnue & 404 & \code{NOT_FOUND}\\
\code{PricingError}, \code{RiskComputationError} & 500 & \code{COMPUTATION_FAILED}\\
Exception inattendue & 500 & \code{INTERNAL_ERROR}\\
\bottomrule
\end{tabularx}
\end{table}

\paragraph{Corrélation par \code{request_id}.} Un middleware ASGI « pur » lit l'en-tête
\code{X-Request-ID} (s'il est bien formé, 1 à 128 caractères sûrs) ou en génère un, le lie à une
\code{ContextVar} et le renvoie dans la réponse. Un filtre de logging l'injecte dans \emph{chaque}
enregistrement, y compris ceux émis par le cœur, qui ignore tout de cette mécanique. Les routes de
calcul sont des fonctions synchrones : FastAPI les exécute dans un \emph{thread}, ce qui évite de
bloquer la boucle d'événements pendant le bootstrap et propage le contexte au thread (ADR-025).

\subsection{Journalisation}

Trois niveaux, au format JSON, un événement par ligne, avec des champs stables : les \emph{logs
techniques} (accès, cache), les \emph{logs de calcul} (par pilier en niveau \code{DEBUG}, synthèse du
bootstrap, PV par jambe) et le \emph{log d'audit} \code{calculation_audit}, qui contient les
empreintes des entrées, les versions et les options, et permet de rejouer le calcul.

\begin{lstlisting}[language=json,caption={Événement de calcul (champs tronqués).}]
{"timestamp": "2026-10-25T00:39:49.930+00:00", "level": "INFO",
 "logger": "core.engines.discounting", "event": "pricing_completed",
 "request_id": "smoke-1", "trade_type": "ois_swap", "currency": "EUR",
 "present_value": 1902.5377349308692, "cashflow_count": 10,
 "par_rate": 0.02254000000000001, "duration_ms": 0.567}
\end{lstlisting}

\subsection{SDK et tableau de bord}

Le \code{PricerClient} (basé sur \code{httpx}) expose \code{price}, \code{build_curves},
\code{catalog}, \code{list_snapshots}, \code{get_snapshot}, \code{health}, \code{ready} et
\code{version}, et renvoie des objets typés qui réutilisent les schémas de l'API. Chaque code d'erreur
devient une exception typée portant le \code{request_id}, à communiquer au support ; une API
injoignable lève \code{ApiUnavailableError}.

\begin{lstlisting}[caption={Utilisation du SDK.}]
from rate_derivatives_pricer.sdk import PricerClient, ois_swap

with PricerClient("http://localhost:8000") as client:
    priced = client.price("EUR_2026-10-19", ois_swap(10_000_000, 0.0225, "pay_fixed", "5Y"))
    print(priced.result.present_value, priced.result.par_rate, priced.request_id)
\end{lstlisting}

Le tableau de bord Streamlit offre deux onglets : la courbe (zéros, forwards, facteurs
d'actualisation, table de calibration) et un formulaire de pricing d'OIS (PV, taux de par,
annuité, PVBP, flux). Il \emph{ne calcule rien} : tout passe par le SDK. Le SDK sert ainsi de test
d'intégration vivant, puisqu'une rupture du contrat HTTP casse immédiatement le tableau de bord.
Le contrat \code{import-linter} (3) interdit à \code{ui} d'importer \code{api} directement ;
le SDK réexporte les schémas de l'API (\code{sdk/models.py}), ce qui rend l'import indirect
\code{ui}~$\to$~\code{sdk}~$\to$~\code{api.schemas} légitime (ADR-021).

% -----------------------------------------------------------------------------
\section{Validation et qualité}\label{sec:validation}

\subsection{QuantLib comme oracle}

QuantLib (version~1.44) n'est utilisée que comme \emph{oracle de test} : dépendance de développement
uniquement, jamais importée par le cœur (ADR-003). Le module \code{quantlib_oracle.py} traduit nos
conventions en objets QuantLib (\code{DepositRateHelper}, \code{OISRateHelper},
\code{PiecewiseLogLinearDiscount}, \code{MakeOIS}). Comme QuantLib possède une date d'évaluation
\emph{globale}, un gestionnaire de contexte la positionne sur la date du snapshot puis restaure la
valeur précédente.

Les cas sont \emph{générés à partir de grilles}, non écrits à la main :
\begin{itemize}
  \item \textbf{Marchés} : 3 dates de valorisation (dont une fin de mois et la veille de Pâques)
        $\times$ 3 chocs parallèles ($-250$, $0$, $+150$ points de base), soit 9 marchés ;
  \item \textbf{Swaps spot} : 13 maturités (1W à 30Y) $\times$ 2 sens $\times$ 2 notionnels
        $\times$ 3 taux fixes $\times$ 9 marchés, soit 1\,404 swaps, plus des swaps déjà commencés
        et à départ différé ;
  \item \textbf{Calendrier} : jours ouvrés TARGET de 1995 à 2060 comparés jour par jour ;
        échéanciers : fréquences $\times$ règles $\times$ fin de mois $\times$ deux ans de dates de
        départ.
\end{itemize}

\subsection{Tolérances et résultats}

Les tolérances sont des constantes nommées d'un seul module de test
(module \code{tolerances.py} du dossier de validation) ; les écarts réellement constatés sont consignés dans le
journal des décisions (ADR-024).

\begin{table}[ht]
\centering
\caption{Écarts maximaux avec QuantLib (9 marchés, 1\,404 swaps spot).}\label{tab:validation}
\begin{tabularx}{\linewidth}{@{}>{\raggedright\arraybackslash}X >{\raggedright\arraybackslash}p{3.2cm} r@{}}
\toprule
Grandeur & Critère & Écart mesuré\\
\midrule
Calendriers, décomptes, échéanciers & égalité exacte & exacte\\
Facteurs d'actualisation, log-linéaire & $<10^{-10}$ & $2{,}5\times10^{-13}$\\
Facteurs d'actualisation, linéaire sur zéro & $<10^{-10}$ & $2{,}3\times10^{-11}$\\
PV rapportée au notionnel & $<10^{-8}$ & $2{,}3\times10^{-13}$\\
Taux de par (relatif) & $<10^{-8}$ & $9{,}9\times10^{-11}$\\
Annuité (relatif) & $<10^{-8}$ & $5{,}6\times10^{-14}$\\
PV au taux de par (propriété) & $\le10^{-10}\times N$ & respecté\\
\bottomrule
\end{tabularx}
\end{table}

La PV est comparée \emph{relativement au notionnel} : un swap proche du par a une PV proche de zéro,
différence de deux jambes de plusieurs millions, pour laquelle un écart relatif à la PV n'a pas de
sens. La figure~\ref{fig:gap} montre l'écart sur les facteurs d'actualisation le long de la courbe de
référence, entre les piliers et au-delà du dernier.

\begin{figure}[ht]
\centering
\begin{tikzpicture}
\begin{axis}[
  width=0.88\linewidth, height=6.4cm, ymode=log, xmin=0, xmax=40, ymin=1e-17, ymax=1e-9,
  xlabel={Maturité (années)}, ylabel={$\lvert P_{\mathrm{pricer}}-P_{\mathrm{QuantLib}}\rvert$},
  legend pos=north west, legend columns=1,
]
\addplot[blue!70!black, thick, only marks, mark=*, mark size=0.8pt] table[x=t, y=loglinear]{data/quantlib_gap.dat};
\addlegendentry{Log-linéaire sur DF (défaut)}
\addplot[orange!85!black, thick, only marks, mark=square*, mark size=0.8pt] table[x=t, y=linearzero]{data/quantlib_gap.dat};
\addlegendentry{Linéaire sur zéro}
\addplot[red!70!black, dashed, thick] coordinates {(0,1e-10) (40,1e-10)};
\addlegendentry{Critère $10^{-10}$}
\end{axis}
\end{tikzpicture}
\caption{Écart absolu sur les facteurs d'actualisation, $0$ à $40$ ans, snapshot de référence. Les écarts
nuls sont tracés au plancher de l'échelle ($10^{-17}$). Maximum : \MaxGapLogLinear{} (log-linéaire) et
\MaxGapLinearZero{} (linéaire sur zéro).}\label{fig:gap}
\end{figure}

\subsection{Pyramide de tests}

\begin{center}
\begin{tabularx}{\linewidth}{@{}l>{\raggedright\arraybackslash}X@{}}
\toprule
Niveau & Contenu\\
\midrule
Unitaires & Calendriers, décomptes, échéanciers, interpolateurs, solveurs, coupons, registres, cache,
mappers ; cas limites et erreurs.\\
Propriétés (\code{hypothesis}) & Sur des centaines d'entrées aléatoires : $\DF{t_0}=1$, facteurs
décroissants, cohérence forward/DF, \textbf{PV nulle au taux de par}, symétrie payeur/receveur,
résidus de bootstrap sous chocs.\\
Validation QuantLib & Comparaison systématique sur les grilles ci-dessus (75 tests).\\
API / SDK / UI & Contrat HTTP, erreurs, \code{request_id} jusque dans les logs du cœur, cache, SDK typé,
tableau de bord exécuté pour de bon (\code{AppTest}).\\
Intégration & SDK~$\to$~HTTP~$\to$~cœur donne \emph{exactement} les mêmes nombres que le cœur appelé
directement ; \emph{golden file} de non-régression.\\
\bottomrule
\end{tabularx}
\end{center}

Au total, \textbf{265 tests}, avec une couverture de \textbf{99\,\%} pour \code{core/} (96\,\% pour
l'ensemble du paquet). Le \emph{golden file} fige la PV, le taux de par et l'annuité de cinq trades
d'exemple ainsi que les facteurs d'actualisation de chaque pilier.

\subsection{Chaîne de qualité}

Avant chaque commit, un hook \code{pre-commit} exécute \code{ruff} (style, docstrings Google,
interdiction de \code{print}), \code{mypy} (mode strict sur \code{core} et \code{services}),
\code{import-linter} et les tests unitaires. La CI GitHub Actions rejoue la même boucle sur une machine
Linux vierge, puis l'ensemble des tests avec mesure de couverture ; la protection de branche exige
qu'elle soit verte avant toute fusion.

\subsection{Incidents rencontrés}

Quatre incidents, instructifs sur la méthode, sont consignés :
\begin{enumerate}
  \item \textbf{Règle de fin de mois.} La première version du générateur déplaçait aussi la date de
        départ ; QuantLib ne le fait pas. \emph{Les détails de convention se vérifient date par date
        contre un oracle, ils ne se devinent pas.}
  \item \textbf{Extrapolation du schéma linéaire sur le zéro.} Linéaire dans la première version, à
        forward constant chez QuantLib ; l'écart n'apparaissait qu'au-delà de 30 ans. \emph{Il faut
        aussi tester hors des piliers.}
  \item \textbf{Contrat \code{import-linter} trop large.} Celui de la phase~0 interdisait à \code{ui}
        d'importer \code{api}, même indirectement, ce qui empêchait le SDK de réutiliser les schémas de
        l'API. Il a été scindé en deux contrats plus précis et la décision est documentée
        (ADR-021).
  \item \textbf{Golden file sous Linux.} Sur la CI, la PV du swap receveur à 30~ans
        ($-1\,034{,}69\eur$, différence de deux jambes d'environ $2{,}7$ millions) s'écartait de
        $2\times10^{-9}\eur$ de la valeur calculée sous Windows, à cause des arrondis d'\code{exp} et
        \code{log} propres à chaque plate-forme : $2{,}2\times10^{-12}$ en relatif à la PV, au-dessus de
        la tolérance relative de $10^{-12}$. Les PV du golden file sont désormais comparées par unité
        de notionnel.
\end{enumerate}

% -----------------------------------------------------------------------------
\section{Journal des décisions}

Les décisions non triviales sont consignées sous forme d'ADR (\emph{Architecture Decision Records}) dans
\code{docs/DECISIONS.md}, dans la même PR que le code concerné. Celles de la phase~1 :

\begin{center}
\small
\begin{tabularx}{\linewidth}{@{}l>{\raggedright\arraybackslash}X@{}}
\toprule
ADR & Décision\\
\midrule
019 & Conventions EUR \estr{} : TARGET, spot lag 2, flux annuels Act/360, Modified Following, fin de
      mois, payment lag 1, sans lookback ; dépôt overnight ; deux variantes 30/360.\\
020 & Fixings : obligatoires avant $t_0$, fixing du jour utilisé s'il est fourni, télescopage de la
      partie projetée, flux échus exclus.\\
021 & Frontières entre couches : schémas partagés API/SDK, entrées du service en dataclasses, contrat
      \code{ui} scindé.\\
022 & Bootstrap : pilier d'un OIS = dernière date de paiement, passe unique, encadrement
      $[-50\,\%,+100\,\%]$, Brent par défaut, extrapolation à forward constant.\\
023 & Données de marché synthétiques.\\
024 & Validation QuantLib : tolérances inchangées, écarts mesurés, PV comparée au notionnel.\\
025 & Contrat HTTP : enveloppe, empreintes, identifiant de requête, erreurs inattendues, variables
      d'environnement, routes synchrones en threads.\\
026 & Clé du cache : SHA-256 des données, de l'interpolation, du solveur et de la version ; nom exclu.\\
\bottomrule
\end{tabularx}
\end{center}

% -----------------------------------------------------------------------------
\section{Limites et suite}\label{sec:limites}

\paragraph{Limites assumées de la phase~1.}
\begin{itemize}
  \item \textbf{Courbe unique.} L'actualisation et la projection partagent la courbe \estr{}. Un swap
        contre Euribor exige une courbe de projection distincte~\cite{ametrano2013,andersen2010}
        (phase~2).
  \item \textbf{Capitalisation simple.} Ni \emph{lookback}, ni \emph{lockout}, ni \emph{observation shift}.
  \item \textbf{Décalage de paiement sans ajustement de convexité} : le coupon est évalué sur la courbe
        puis actualisé à sa date de paiement ; l'écart est cohérent avec le moteur d'actualisation de
        QuantLib, sur lequel la validation s'appuie.
  \item \textbf{Une seule devise} (EUR) ; l'USD (SOFR, calendrier US) s'ajoutera par configuration.
  \item \textbf{Données synthétiques} : le snapshot peut être remplacé par des cotations réelles sans
        changement de code.
  \item \textbf{Pas de risque} : seuls le PVBP et l'annuité sont fournis.
\end{itemize}

\paragraph{Prochaines étapes.} La phase~2 introduit l'\code{IborIndex} (Euribor 3M/6M), le bootstrap
multi-courbes (actualisation OIS, projection Euribor), le swap vanille (IRS) et le risque par
\emph{bump-and-reprice} sur les instruments de calibration : DV01 parallèle, DV01 par bucket
(\emph{key-rate}), convexité et DV01 par courbe. L'architecture a été conçue pour cela : le risque
dépendra d'une fonction de valorisation abstraite, de sorte qu'ajouter une mesure ne demandera
aucune modification des moteurs.

% =============================================================================
\appendix
\section{Calibration détaillée de la courbe}\label{app:pillars}

\begin{table}[H]
\centering
\caption{Piliers calibrés sur le snapshot synthétique du \ValuationDate{} : cotation, facteur
d'actualisation, taux zéro continu, itérations de Brent et résidu final (taux décimal).}
\setlength{\tabcolsep}{4pt}
\small
\begin{tabular}{@{}llcrrrcr@{}}
\toprule
Pilier & Type & Date & Cotation (\%) & DF & Zéro (\%) & Itér. & Résidu \\
\midrule
ON & Dépôt & 2026-10-20 & 1{,}921 & 0{,}99994664 & 1{,}948 & 5 & $-6{,}0\times10^{-15}$ \\
1W & OIS & 2026-10-29 & 1{,}922 & 0{,}99946637 & 1{,}948 & 6 & $-8{,}4\times10^{-16}$ \\
2W & OIS & 2026-11-05 & 1{,}921 & 0{,}99909360 & 1{,}947 & 6 & $-3{,}4\times10^{-16}$ \\
1M & OIS & 2026-11-24 & 1{,}917 & 0{,}99808630 & 1{,}942 & 7 & $4{,}8\times10^{-16}$ \\
2M & OIS & 2026-12-22 & 1{,}905 & 0{,}99662389 & 1{,}929 & 7 & $-1{,}3\times10^{-15}$ \\
3M & OIS & 2027-01-22 & 1{,}893 & 0{,}99502796 & 1{,}915 & 7 & $2{,}9\times10^{-16}$ \\
6M & OIS & 2027-04-22 & 1{,}874 & 0{,}99045842 & 1{,}892 & 7 & $-7{,}4\times10^{-16}$ \\
9M & OIS & 2027-07-22 & 1{,}871 & 0{,}98585436 & 1{,}884 & 7 & $4{,}9\times10^{-16}$ \\
1Y & OIS & 2027-10-22 & 1{,}885 & 0{,}98109021 & 1{,}894 & 7 & $-3{,}0\times10^{-16}$ \\
18M & OIS & 2028-04-24 & 1{,}934 & 0{,}97095063 & 1{,}946 & 8 & $-5{,}2\times10^{-17}$ \\
2Y & OIS & 2028-10-24 & 1{,}992 & 0{,}96046173 & 2{,}001 & 7 & $-1{,}4\times10^{-16}$ \\
3Y & OIS & 2029-10-23 & 2{,}098 & 0{,}93845492 & 2{,}108 & 8 & $-6{,}9\times10^{-18}$ \\
4Y & OIS & 2030-10-22 & 2{,}183 & 0{,}91575196 & 2{,}194 & 7 & $7{,}6\times10^{-17}$ \\
5Y & OIS & 2031-10-22 & 2{,}254 & 0{,}89261592 & 2{,}267 & 7 & $1{,}4\times10^{-17}$ \\
6Y & OIS & 2032-10-22 & 2{,}318 & 0{,}86908808 & 2{,}333 & 7 & $-3{,}5\times10^{-18}$ \\
7Y & OIS & 2033-10-24 & 2{,}379 & 0{,}84512772 & 2{,}397 & 7 & $-1{,}4\times10^{-17}$ \\
8Y & OIS & 2034-10-24 & 2{,}436 & 0{,}82115870 & 2{,}457 & 7 & $2{,}1\times10^{-17}$ \\
9Y & OIS & 2035-10-23 & 2{,}488 & 0{,}79728635 & 2{,}513 & 7 & $-2{,}1\times10^{-17}$ \\
10Y & OIS & 2036-10-22 & 2{,}535 & 0{,}77357667 & 2{,}563 & 7 & $-6{,}9\times10^{-18}$ \\
12Y & OIS & 2038-10-22 & 2{,}612 & 0{,}72756799 & 2{,}647 & 7 & $3{,}5\times10^{-18}$ \\
15Y & OIS & 2041-10-22 & 2{,}688 & 0{,}66362996 & 2{,}730 & 7 & $0$ \\
20Y & OIS & 2046-10-23 & 2{,}731 & 0{,}57388463 & 2{,}773 & 9 & $-1{,}0\times10^{-17}$ \\
25Y & OIS & 2051-10-24 & 2{,}702 & 0{,}50544470 & 2{,}726 & 8 & $-3{,}5\times10^{-18}$ \\
30Y & OIS & 2056-10-24 & 2{,}651 & 0{,}45122111 & 2{,}650 & 8 & $-1{,}4\times10^{-17}$ \\
\bottomrule
\end{tabular}

\end{table}

\section{Carte des modules}\label{app:modules}

Chemins relatifs à \code{src/rate_derivatives_pricer/}.

\begingroup
\small
\begin{longtable}{@{}>{\raggedright\arraybackslash}p{6.5cm}>{\raggedright\arraybackslash}p{8.9cm}@{}}
\toprule
Module & Rôle\\
\midrule
\endhead
\code{core/errors.py} & Hiérarchie d'exceptions typées.\\
\code{core/registry/} & Registre générique et six registres concrets, chargement paresseux.\\
\code{core/conventions/business_day.py} & Conventions de jour ouvré.\\
\code{core/conventions/tenor.py} & \code{Tenor}, \code{Frequency}, arithmétique de dates.\\
\code{core/conventions/calendar.py} & Calendriers TARGET, week-ends, nul, joint ; fin de mois.\\
\code{core/conventions/day_count.py} & Act/360, Act/365F, 30/360, 30E/360.\\
\code{core/conventions/currency.py} \newline \code{core/conventions/data/EUR.json} &
  Conventions d'une devise, chargées depuis la configuration.\\
\code{core/schedule/generator.py} & Génération d'échéanciers (accrual, paiement, stub).\\
\code{core/market/indices.py} & Description d'un indice overnight.\\
\code{core/market/quotes.py} & Cotations avec unité explicite.\\
\code{core/market/snapshot.py} & Snapshot immuable : date, cotations, fixings.\\
\code{core/market/providers.py} & Fournisseur de fichiers, format canonique.\\
\code{core/market/interpolation.py} & Interpolations log-linéaire sur DF et linéaire sur zéro.\\
\code{core/market/solvers.py} & Solveurs de Brent et de Newton.\\
\code{core/market/curves.py} & Courbe d'actualisation (DF, zéro, forward).\\
\code{core/market/environment.py} & Environnement de marché, rapports de calibration.\\
\code{core/market/calibration.py} & Instruments de calibration (dépôt, OIS).\\
\code{core/market/bootstrap.py} & \code{CurveBuilder} : bootstrap séquentiel.\\
\code{core/instruments/cashflows.py} & Flux polymorphes : coupon fixe, coupon overnight composé.\\
\code{core/instruments/legs.py} & Jambe : flux, notionnel, sens.\\
\code{core/instruments/swaps.py} & \code{OisSwap}, \code{SwapDirection}.\\
\code{core/instruments/factories.py} & Dates spot et maturité, construction du swap, factory.\\
\code{core/engines/base.py} & Interface de moteur, sélection par registre.\\
\code{core/engines/discounting.py} & Moteur d'actualisation d'un swap.\\
\code{core/engines/results.py} & \code{PricingResult} et sous-résultats.\\
\code{services/requests.py} & Entrées du service (dataclasses figées).\\
\code{services/mappers.py} & Traductions, empreintes, charges utiles JSON.\\
\code{services/curve_cache.py} & Cache LRU de courbes, protocole \code{CurveCache}.\\
\code{services/pricer_service.py} & Façade : calculs, métadonnées, diagnostics, audit.\\
\code{api/settings.py} & Configuration par variables d'environnement.\\
\code{api/schemas.py} & Contrat HTTP (Pydantic).\\
\code{api/errors.py} & Traduction des erreurs en réponses structurées.\\
\code{api/middleware.py} & \code{request_id}, journal d'accès, erreurs inattendues.\\
\code{api/routes.py} & Routes de la version~1.\\
\code{api/main.py} & Fabrique de l'application.\\
\code{api/logging_config.py} & Logging JSON et \code{request_id} (phase~0).\\
\code{sdk/client.py}, \code{sdk/models.py} & Client typé et modèles réexportés.\\
\code{ui/views.py}, \code{ui/app.py} & Présentation (figures, tables) et script Streamlit.\\
\bottomrule
\end{longtable}
\endgroup

\section{Reproduire les chiffres et compiler}

Le script \code{docs/latex/generate_data.py} appelle le pricer et écrit, à côté de ce document, les
données des figures (\code{data/*.dat}), les tableaux (\code{generated/table_*.tex}) et les macros des
exemples chiffrés (\code{generated/numbers.tex}). Pour régénérer puis compiler :

\begin{lstlisting}[language={},caption={Régénération des données et compilation.}]
uv run python docs/latex/generate_data.py
cd docs/latex
pdflatex pricing_ois.tex
pdflatex pricing_ois.tex      # second passage : table des matieres et references
\end{lstlisting}

% =============================================================================
\begin{thebibliography}{9}

\bibitem{andersen2010}
L.~B.~G.~Andersen et V.~V.~Piterbarg, \emph{Interest Rate Modeling}, vol.~1, Atlantic Financial Press, 2010.

\bibitem{ametrano2013}
F.~M.~Ametrano et M.~Bianchetti, \emph{Everything You Always Wanted to Know About Multiple Interest
Rate Curve Bootstrapping but Were Afraid to Ask}, SSRN, 2013.

\bibitem{brent1973}
R.~P.~Brent, \emph{Algorithms for Minimization without Derivatives}, Prentice-Hall, 1973 (chapitre~4).

\bibitem{hagan2006}
P.~S.~Hagan et G.~West, \emph{Interpolation Methods for Curve Construction}, Applied Mathematical
Finance, 13(2), 2006.

\bibitem{henrard2014}
M.~Henrard, \emph{Interest Rate Modelling in the Multi-Curve Framework}, Palgrave Macmillan, 2014.

\bibitem{isda2006}
ISDA, \emph{2006 ISDA Definitions}, section~4.16 (conventions de décompte) et section~4.12
(conventions de jour ouvré).

\bibitem{piterbarg2010}
V.~V.~Piterbarg, \emph{Funding Beyond Discounting: Collateral Agreements and Derivatives Pricing},
Risk, février 2010.

\end{thebibliography}

\end{document}
